% ============================================================
% 字体配置文件
% 使用 XeLaTeX 引擎，通过 fontspec 设置中英文字体
% ============================================================
\usepackage{fontspec}
\usepackage{xeCJK}
% ---------- 英文字体（衬线 / 无衬线 / 等宽） ----------
\setmainfont{Times New Roman}
\setsansfont{Arial}
\setmonofont{Courier New}

% ---------- 中文字体（ctex 默认字体集） ----------
% ctexart 自动加载中文字体，无需额外设置
% 如需自定义，取消下面的注释并修改字体名：
% \setCJKsansfont{SimHei}
% \setCJKmonofont{FangSong}
% ===== 中文正文（默认）=====
\setCJKmainfont{SimSun}[BoldFont=SimHei, ItalicFont=KaiTi]
%中文字库通常没有独立的“粗体”或“斜体”字形文件，所以这里通过 BoldFont 和 ItalicFont 参数指定了映射关系

% ===== 中文无衬线（标题）=====
\setCJKsansfont{SimHei}

% ---------- 数学字体 ----------
% 如需自定义数学字体，需先在 packages.tex 中加载 unicode-math 宏包，
% 然后取消下面的注释：
% \setmathfont{Latin Modern Math}
